% Sección D · Métrica y sensibilidad.
% El desarrollo (datos de entrada, cálculo tier por tier, métodos alternativos y
% las otras tres decisiones) está en el anexo D, para respetar el límite de 6 páginas.

\section{Métrica y sensibilidad}
\label{sec:metrica}

\paragraph{Qué medimos.} La métrica es el precio medio por litro de exportación de cada semestre de 2025:
\[
P = \frac{\sum_i M_i}{\sum_i L_i}, \qquad L_i = C_i \times k_i,
\]
donde, para cada línea de venta $i$, $M_i$ es el monto neto en USD (después de descuento), $L_i$ son los litros vendidos, $C_i$ son las cajas y $k_i$ son los litros por caja de su producto. La suma recorre las líneas de \texttt{fact\_venta} del canal exportación con fecha de factura en el semestre (10.672 en el año), incluidas las de merma, con signo negativo. Es una razón de sumas, no un promedio de precios. Entre el 1.\textsuperscript{er} y el 2.\textsuperscript{o} semestre, $P$ subió de 4,574 a 5,085 USD/L: +0,511 USD/L, un 11,2\,\%.

\paragraph{Por qué no basta con mirar $P$.} Un precio medio puede subir por dos razones distintas: porque cada vino se vende más caro (\emph{efecto precio}) o porque, con los mismos precios, se vende una mayor proporción del vino caro (\emph{efecto mezcla}). Para separarlas usamos la descomposición estándar precio-mezcla (Kitagawa, 1955; Horngren, Datar y Rajan, \emph{Cost Accounting}): para cada tier $t$, con $w_t$ su proporción de los litros y $p_t$ su precio por litro, el cambio de $P$ se reparte en tres términos que suman exactamente el alza: mezcla $\sum_t (w_{t,2}-w_{t,1})\,p_{t,1}$, precio $\sum_t w_{t,1}\,(p_{t,2}-p_{t,1})$ y cruzado $\sum_t (w_{t,2}-w_{t,1})(p_{t,2}-p_{t,1})$. El cálculo tier por tier está en el anexo~D. El criterio, fijado antes de mirar los resultados: el alza es de mezcla si ese efecto explica más de la mitad del cambio con todos los métodos y definiciones alternativas. No usamos pruebas de significancia: los datos son todas las ventas del año, no una muestra, y la incertidumbre viene de nuestras definiciones, que es lo que miden las sensibilidades.

\paragraph{Resultado.}

\begin{center}
\small
\begin{tabular}{l r r}
\hline
\textbf{Efecto} & \textbf{USD/L} & \textbf{\% del alza} \\
\hline
Mezcla & +0,551 & +108\,\% \\
Precio & $-$0,035 & $-$7\,\% \\
Cruzado & $-$0,005 & $-$1\,\% \\
\hline
\textbf{Cambio total} & \textbf{+0,511} & \textbf{100\,\%} \\
\hline
\end{tabular}
\end{center}

\noindent Vender 9,8 puntos más de gran reserva, el tier más caro, suma 0,81 USD/L al precio medio; vender 9,3 puntos menos de varietal resta 0,23. El efecto precio es negativo en los tres tiers: los precios de lista bajaron entre 0,15\,\% y 0,19\,\% y el descuento medio subió de 3,82\,\% a 4,40\,\% (ponderados por cajas). El cruzado es el 1\,\% del alza, así que la repartición no es ambigua. Según el criterio fijado, \textbf{el alza es de mezcla}.

\paragraph{La decisión de definición: moneda y fecha de conversión.}\label{par:fecha_dolar} La viña vende en dólares (V), pero sus costos y su caja están en pesos, y finanzas mira el resultado en pesos. Convertir exige elegir con qué dólar: el del día de la factura o el del día en que se cobra (factura más 60 días en el hemisferio sur y 90 en el norte, según la contraparte (V)). Cada línea se convierte con el dólar observado de su propia fecha, nunca con un promedio del período. Esta decisión sí mueve el número:

\begin{center}
\small
\renewcommand{\arraystretch}{1.15}
\begin{tabular}{@{}l r@{}}
\toprule
\textbf{Definición} & \textbf{Alza del precio medio por litro} \\
\midrule
USD (adoptada) & +11,2\,\% \\
CLP, dólar de la fecha de factura & +10,1\,\% \\
CLP, dólar de la fecha de cobranza & +6,3\,\% \\
\bottomrule
\end{tabular}
\end{center}

\noindent Según la fecha que se elija, el alza en pesos va de 10,1\,\% a 6,3\,\%: casi la mitad del 11,2\,\% en dólares. La diferencia es el dólar, que bajó entre semestres: los mismos dólares se convirtieron en menos pesos, y con la fecha de cobranza la caída es mayor porque el dólar se toma aún más tarde (M). Descartamos la fecha de embarque: el 30\,\% de las líneas de embarque tiene fecha anterior a la factura (M).

\paragraph{Qué cambia, qué no y por qué nos quedamos con dólares.} Cambia el tamaño del alza en pesos, no su causa: la caída del dólar no es ni precio ni mezcla, solo le resta al alza (I). Medimos en dólares porque los contratos se firman en esa moneda (V) y así la mezcla queda separada del tipo de cambio. La versión en pesos, la que llega a caja y mira finanzas, se elige en el prototipo.

\paragraph{Las demás decisiones no cambian la conclusión.} Para cada una repetimos el análisis con la opción que no elegimos:

\begin{center}
\small
\renewcommand{\arraystretch}{1.15}
\begin{tabular}{@{}l l l r@{}}
\toprule
\textbf{Decisión} & \textbf{Adoptada} & \textbf{Alternativa} & \textbf{Mezcla, \% del alza} \\
\midrule
Litros por caja magnum & 9,0 L (maestro de costos) & 18 o 4,5 L & 108,2\,\% \\
Correcciones de limpieza & Aplicadas & Sin aplicar (ERP) & 78,2\,\% \\
\bottomrule
\end{tabular}
\label{tab:sensibilidades}
\end{center}

\noindent Con la definición adoptada, la mezcla explica el 107,9\,\% del alza, y en ambas alternativas supera el 50\,\% del criterio. Sin aplicar las correcciones baja a 78,2\,\%: las cajas cargadas en botellas y los montos mal registrados meten ruido en el precio por litro, pero aun así la mezcla explica la mayor parte del alza. Las alternativas de 18 y 4,5\,L por caja son extremos hipotéticos que no cambian la respuesta; su detalle y las demás sensibilidades de la Entrega~1 (método de descomposición, merma y precio de lista) están en el anexo~D.